% Generated draft: requires mathematical and source review.
\documentclass[11pt]{article}
\usepackage[margin=25mm]{geometry}
\usepackage{fontspec}
\setmainfont{Latin Modern Roman}
\usepackage{amsmath,amssymb,mathtools,unicode-math}
\setmathfont{Latin Modern Math}
\usepackage{enumitem,needspace,xurl,hyperref}
\setlength{\parindent}{0pt}
\setlength{\parskip}{5pt}
\setlength{\emergencystretch}{3em}
\newcommand{\N}{\mathbb{N}}
\newcommand{\Z}{\mathbb{Z}}
\newcommand{\Q}{\mathbb{Q}}
\newcommand{\R}{\mathbb{R}}
\newcommand{\C}{\mathbb{C}}
\newcommand{\F}{\mathbb{F}}
\newcommand{\A}{\mathbb{A}}
\newcommand{\PP}{\mathbb{P}}
\newcommand{\E}{\mathbb{E}}
\newcommand{\eps}{\varepsilon}
\newcommand{\dd}{\,\mathrm{d}}
\begin{document}
\section*{UnsolvedMath ID: 30006553}
\textbf{Source problem code:} catalogue entry
\subsection{Definitions and mathematical statement}
The supplied question concerns a local version of graph removal. To describe the distance under a conventional graph model, let \(G=(V,E)\) and \(H\) be finite simple undirected graphs, with \(n=|V|\geq 1\). These graph conventions are provisional: the supplied record does not explicitly specify them. Say that \(G\) is \(H\)-free if it contains no subgraph isomorphic to \(H\); the source does not explicitly distinguish ordinary from induced copies.

For a set \(F\subseteq E\) of deleted edges, define its removed-edge degree at \(v\in V\) by
\[
 d_F(v)=|\{e\in F:v\in e\}|,
 \qquad
 \|d_F\|_\infty=\max_{v\in V}d_F(v).
\]
Under the ordinary-subgraph interpretation, the corresponding deletion distance is
\[
 D_{\infty,H}(G)=
 \min\bigl\{\|d_F\|_\infty:
 F\subseteq E,\ (V,E\setminus F)\text{ is }H\text{-free}\bigr\}.
\]
Its vertex-normalized version is \(D_{\infty,H}(G)/n\). The source does not specify whether the distance is normalized or which normalization is intended.

\textbf{Source question.} Does an analogue of the graph removal lemma hold when distance from being \(H\)-free is measured by the maximum number of deleted edges incident to any one vertex, rather than by the total number of deleted edges?

\textbf{Unresolved formulation.} The supplied record does not give the copy-counting hypothesis, its normalization, the order and dependencies of the quantitative parameters, or the range of graphs \(H\). In particular, it does not specify whether the hypothesis bounds the total number of copies of \(H\), bounds copies through individual vertices, or imposes another local condition. Therefore the displayed distance defines a provisional interpretation of the metric, but the evidence does not determine a unique fully quantified removal conjecture. No particular copy-counting implication is asserted here.
\subsection{Short English statement}
Can graph removal be formulated so that making a graph \(H\)-free requires deleting only a small number of edges at each vertex? The supplied record leaves the necessary copy-counting hypothesis and quantitative formulation unspecified, so a precise complete conjecture requires clarification from the source.
\subsection{Sources}
\begin{itemize}
\item Supplied local UnsolvedMath registry record, permanent ID \(30006553\), title: \emph{Graph Removal in the \(L_\infty\) Metric}; exact source problem code: \texttt{catalogue entry}. The record belongs to the Oberwolfach Reports Open Problems collection and gives proposed year 2026.
\item Supplied report citation: \emph{Combinatorics} (2026), Oberwolfach Reports, DOI 10.4171/owr/2026/1, pp.~5--79. Source URL: \url{https://doi.org/10.4171/owr/2026/1}. The same URL is supplied in the literature sources under the label ``Oberwolfach Report 1/2026.'' No problem-specific page or full quantitative statement is supplied.
\item Attribution in the original statement: ``Local Removal Lemmas,'' Fan Wei, joint work with Jacob Fox, Hung Le, and Ruilin Shi. The registry does not separately identify a proposer.
\item Literature assessment date: \texttt{2026-08-22}. This is supplied dated curation, not independently verified. The supplied status is ``partially solved'': the assessment reports results for cliques and several other families, including bipartite and odd-cycle cases, while describing the general graph statement as conjectural. The precise formulations and hypotheses of these reported results are not supplied. The background explicitly requests independent verification of time-sensitive status and cited claims.
\item The registry labels the statement \texttt{corrected\_verified} and reports that the official report was checked and notation repaired. These are supplied provenance claims; no independent source verification is asserted in this draft.
\end{itemize}
\end{document}
